\documentclass{article}

\usepackage{preamb}

\usepackage[english]{babel}
\usepackage[utf8]{inputenc}
\usepackage[backend=biber,style=verbose,citetracker=true]{biblatex}
\usepackage[T1]{fontenc}
\usepackage{csquotes}

\graphicspath{{../images/}}
\addbibresource{bibliography.bib}

\author{Tanguy Lefort}
\title{Title}


%%%%%%%%%%%%%%%%%%%%%%%%%%%%%%%%%%%%%%%%%%%%%%%%%%%%%%%%%%%%%%%%%%%%
%%%%%%%%%%%%%%%%%%%% Begin document - title page %%%%%%%%%%%%%%%%%%%
%%%%%%%%%%%%%%%%%%%%%%%%%%%%%%%%%%%%%%%%%%%%%%%%%%%%%%%%%%%%%%%%%%%%


\begin{document}

\begin{titlepage} 
	\newcommand{\HRule}{\rule{\linewidth}{0.5mm}}
	
	\center
	
	\textsc{\LARGE University of Montpellier}\\[1.5cm]
	
	\textsc{\Large Master 2 Biostatistics }\\[0.5cm] 
	
	\textsc{\large HMMA$307$ project}\\[0.5cm] 
    \vspace{2cm}
	\HRule\\[0.4cm]
	
	{\huge\bfseries Sparse Regularization via Bidualization}\\[0.4cm] 
	
	\HRule\\[1.5cm]
	
	\begin{minipage}{0.4\textwidth}
		\begin{flushleft}
			\large
			\textit{By:}\\
			Tanguy Lefort 
		\end{flushleft}
	\end{minipage}
	~
	\begin{minipage}{0.4\textwidth}
		\begin{flushright}
			\large
			\textit{Teacher: }\\
			Joseph salmon\\
		\end{flushright}
	\end{minipage}
	\vspace{2cm}
	\vfill % Position the date 3/4 down the remaining page
	
	{\large 2020} 
	\begin{tikzpicture}[remember picture,overlay]
    \node[anchor=south] at (current page text area.south) {\includegraphics[scale=0.3]{Logo.pdf}};
    \end{tikzpicture}	
	 
	%----------------------------------------------------------------------------------------
	
	\vfill 
\end{titlepage}

\tableofcontents

%%%%%%%%%%%%%%%%%%%%%%%%%%%%%%%%%%%%%%%%%%%%%%%%%%%%%%%%%%%%%%%%%%%%
% Introduction
%%%%%%%%%%%%%%%%%%%%%%%%%%%%%%%%%%%%%%%%%%%%%%%%%%%%%%%%%%%%%%%%%%%%

\section*{Introduction}

The increasing................


%%%%%%%%%%%%%%%%%%%%%%%%%%%%%%%%%%%%%%%%%%%%%%%%%%%%%%%%%%%%%%%%%%%%
% Sections of presentation
%%%%%%%%%%%%%%%%%%%%%%%%%%%%%%%%%%%%%%%%%%%%%%%%%%%%%%%%%%%%%%%%%%%%

\section{Solving problems with regularization}

% Present the problem with existing ridge and why L0
Regularized problems are a very useful way to obtain a solution in many situations. We write the problem as 
\[\min_{x\in\RR^n} f(x) + \mathrm{pen}(x)\enspace,\]
where $f$ only depends of the data collected and $\mathrm{pen}$ is a penalty term. It can represent for example a constraint over the number of variables that play a role.
In the case of the Tikhonov-ridge regularization, for a matrix $A\in\RR^{n\times n}$ made from the data and an observed signal $b\in\RR^n$, we write the problem
\[\min_{x\in\RR^n} \frac{1}{2}\|Ax-b\|^2_2 + \frac{\lambda}{2}\|x\|^2_2\enspace,\]
where $\lambda \geq 0$ represents how strong the penalty on the $l_2$ norm of the solution is. Note that for $\lambda = 0$ we retrieve the least square problem, and when $\lambda\rightarrow +\infty$, we have $x=0$.
However, when working with an observed signal, we sometimes know in advance the sparsity (or a reasonable estimate) of the original one. So we have an information that is stronger than a constraint on the norm of the vector, we have a constraint on the number of non-zero values, meaning on the $l_0$ norm.

% define the l_0 norm and space Ck
\begin{definition}
Let $x\in\RR^n$, the $l_0$ norm of $x$ noted $\|x\|_0$ is the number of non-zero values: \[\|x\|_0:=\abs{\{i\,|\, x_i\neq 0\}}\enspace .\]
We note $C_k$ the space of the $k$-sparse vectors: $C_k := \{x\in\RR^n\,|\, \|x\|_0 = k\}.$
\end{definition}

% introduce LASSO and ENET
One of the main advantages of the Tikhonov ($l_2$) regularization is the differentiability, the convexity and the continuity. With the $l_0$ norm, we lose them all. This lead to a compromise called the \lasso using the $l_1$ norm to keep the continuity and use sub-differentials instead of regular ones. The \lasso method will select some variables (depending on the value of $\lambda$) to explain the observed signal. It is very useful in a lot of situations like in genomics, but in situations with highly correlated groups of variables, the \lasso tends to only select one from each group and force the others to a null-effect. This issue lead to the development of the \enet method which uses a penalizer that is a combination of the Tikhonov and \lasso penalty
\[\min_{x\in\RR^n} \frac{1}{2}\|Ax-b\|^2_2 + \lambda_1 \|x_1\| +  \frac{\lambda_2}{2}\|x\|^2_2\enspace.\]

% the studied penalty
The proposition of \textbf{MAKE CITATION} is to keep the $l_0$ norm for the sparsity penalty and use the $l_2$ norm in the \enet for the grouping effect. The penalty is then:
\[s_k(x)=\begin{cases} \frac{1}{2}\|x\|_2^2, \text{ if } \|x\|_0\leq k\\
\infty, \text{ else}\end{cases} = \frac{1}{2}\|x\|_2^2 + \delta_{C_k}\enspace.\]
However, the first problem is still there. We don't have neither the continuity nor the convexity. That's why instead of considering $s_k$, we will consider it's best convex estimator, the sparse envelope:
\[ \mathcal{S}_k(x) := s_k^{**}(x)\enspace,\]
where, for a function $f$, we define its convex conjugate by  $f^*(y)=\max_{x\in\RR^n}\{\langle x, y\rangle - g(x)\}.$
Our goal from now is to compute efficiently the proximal operator of $\lambda \mathcal{S}_k$, for $\lambda >0$ and then compare our results to the \enet ones.


%%%%%%%%%%%%%%%%%%%%%%%%%%%%%%%%%%%%%%%%%%%%%%%%%%%%%%%%%%%%%%%%%%%%
% Sections 2: s** = S (what is this object?)
%%%%%%%%%%%%%%%%%%%%%%%%%%%%%%%%%%%%%%%%%%%%%%%%%%%%%%%%%%%%%%%%%%%%
\section{Compute the convex biregularization}

%%%%%%%%%%%%%%%%%%%%%%%%%%%%%%%%%%%%%%%%%%%%%%%%%%%%%%%%%%%%%%%%%%%%
% Sections 3: prox_lambdaSk (comparisons with Enet)
%%%%%%%%%%%%%%%%%%%%%%%%%%%%%%%%%%%%%%%%%%%%%%%%%%%%%%%%%%%%%%%%%%%%
\section{Proximal operator computation and numerical results}



%%%%%%%%%%%%%%%%%%%%%%%%%%%%%%%%%%%%%%%%%%%%%%%%%%%%%%%%%%%%%%%%%%%%
% Conclusion
%%%%%%%%%%%%%%%%%%%%%%%%%%%%%%%%%%%%%%%%%%%%%%%%%%%%%%%%%%%%%%%%%%%%

\section*{Conclusion}

In conclusion, the biregularization method................


%%%%%%%%%%%%%%%%%%%%%%%%%%%%%%%%%%%%%%%%%%%%%%%%%%%%%%%%%%%%%%%%%%%%%%%%%%%%%
%%%%%%%%%%%%%%%%%%%%%%%%%%    Bibliography   %%%%%%%%%%%%%%%%%%%%%%%%%%%%%%%%
%%%%%%%%%%%%%%%%%%%%%%%%%%%%%%%%%%%%%%%%%%%%%%%%%%%%%%%%%%%%%%%%%%%%%%%%%%%%%

\section{Bibliography}
\nocite{*}
\printbibliography


\end{document}
